%==============================================================================
%  CAPSTONE PROJECT: Optimizer Showdown
%  Part II: Optimization Theory
%==============================================================================
\documentclass[../../main.tex]{subfiles}

\begin{document}

\chapter{Capstone Project: Optimizer Showdown}
\label{ch:capstone-optimizers}

\begin{keytakeaways}
In this project, you will:
\begin{itemize}
    \item Implement SGD, Momentum, RMSprop, and Adam from scratch
    \item Train a neural network on MNIST using each optimizer
    \item Compare convergence speed, stability, and final accuracy
    \item Understand when to use each optimizer in practice
\end{itemize}
Estimated time: 3--4 hours
\end{keytakeaways}

\begin{realworld}
\textbf{Why This Matters:} Understanding optimizers deeply---not just using them as black boxes---separates good practitioners from great ones. When training fails, knowing \textit{why} Adam oscillates or SGD stalls lets you diagnose and fix problems instead of blindly trying hyperparameters.

\textbf{What You'll Build:} A complete optimizer benchmarking framework that you can reuse for any future project. The patterns here scale from MNIST to million-parameter models.
\end{realworld}

%------------------------------------------------------------------------------
\section{The Challenge}
\label{sec:capstone-challenge}

Your mission: train a simple multi-layer perceptron (MLP) on the MNIST handwritten digit dataset using four different optimizers, and systematically compare their performance.

\subsection{The Network Architecture}

We will use a straightforward three-layer MLP:
\begin{itemize}
    \item \textbf{Input layer}: 784 neurons (28$\times$28 flattened images)
    \item \textbf{Hidden layer 1}: 256 neurons with ReLU activation
    \item \textbf{Hidden layer 2}: 128 neurons with ReLU activation
    \item \textbf{Output layer}: 10 neurons with softmax activation
\end{itemize}

This architecture has approximately 235,000 trainable parameters---large enough to see meaningful optimizer differences, small enough to train quickly on a laptop.

\subsection{The Competition Criteria}

We will compare the four optimizers on:

\begin{enumerate}
    \item \textbf{Convergence speed}: How many epochs to reach 95\% validation accuracy?
    \item \textbf{Final accuracy}: What is the best validation accuracy achieved during training?
    \item \textbf{Training stability}: How smooth is the loss curve? Are there oscillations?
    \item \textbf{Learning rate sensitivity}: How does performance change with different learning rates?
\end{enumerate}

\begin{important}
\textbf{Proper Evaluation Protocol}

We use a three-way data split to ensure unbiased evaluation:
\begin{itemize}
    \item \textbf{Training set (80\%)}: Used to update model parameters
    \item \textbf{Validation set (10\%)}: Used to compare optimizers during training and for early stopping decisions
    \item \textbf{Test set (10\%)}: Evaluated only once at the very end for final reporting
\end{itemize}

Using the test set during training would leak information into model selection, giving overly optimistic estimates of generalization performance. This is a common mistake---even in published papers!
\end{important}

\begin{keyconcept}
Before starting to code, write down your predictions:
\begin{itemize}
    \item Which optimizer do you expect to converge fastest?
    \item Which will achieve the highest final accuracy?
    \item Which will be most sensitive to learning rate choice?
\end{itemize}
Comparing your predictions to actual results deepens understanding.
\end{keyconcept}

%------------------------------------------------------------------------------
\section{Quick Reference: Key Formulas}
\label{sec:capstone-formulas}

Before implementing, here is a quick refresher of each optimizer's update rule. For detailed explanations, refer to Chapter~\ref{ch:gradient-descent} (momentum) and Chapter~\ref{ch:advanced-optimizers} (Adam, schedules).

\begin{table}[htbp]
\centering
\caption{Optimizer formulas used in this project}
\label{tab:optimizer-formulas}
\renewcommand{\arraystretch}{1.4}
\small
\begin{tabular}{@{}lp{10cm}@{}}
\toprule
\textbf{Method} & \textbf{Update Rule} \\
\midrule
SGD & $\vect{\theta}_{t+1} = \vect{\theta}_t - \eta \nabla L(\vect{\theta}_t)$ \\[8pt]
Momentum & $\vect{v}_{t+1} = \beta \vect{v}_t + \nabla L$; \quad $\vect{\theta}_{t+1} = \vect{\theta}_t - \eta \vect{v}_{t+1}$ \\[8pt]
RMSprop & $\vect{s}_{t+1} = \rho \vect{s}_t + (1-\rho)g_t^2$; \quad $\displaystyle\vect{\theta}_{t+1} = \vect{\theta}_t - \frac{\eta}{\sqrt{\vect{s}_{t+1}+\epsilon}} \odot g_t$ \\[12pt]
Adam & \begin{tabular}[t]{@{}l@{}}
$\vect{m}_t = \beta_1 \vect{m}_{t-1} + (1-\beta_1)g_t$ \\[4pt]
$\vect{v}_t = \beta_2 \vect{v}_{t-1} + (1-\beta_2)g_t^2$ \\[4pt]
$\displaystyle\vect{\theta}_{t+1} = \vect{\theta}_t - \eta \frac{\hat{\vect{m}}_t}{\sqrt{\hat{\vect{v}}_t}+\epsilon}$
\end{tabular} \\[6pt]
\bottomrule
\end{tabular}
\end{table}

\begin{table}[htbp]
\centering
\caption{Typical hyperparameter values}
\label{tab:optimizer-hyperparams}
\begin{tabular}{@{}llll@{}}
\toprule
\textbf{Optimizer} & \textbf{Learning Rate} & \textbf{Momentum/Decay} & \textbf{Other} \\
\midrule
SGD & 0.1 & --- & --- \\
Momentum & 0.01 & $\beta = 0.9$ & --- \\
RMSprop & 0.001 & $\rho = 0.99$ & $\epsilon = 10^{-8}$ \\
Adam & 0.001 & $\beta_1 = 0.9$, $\beta_2 = 0.999$ & $\epsilon = 10^{-8}$ \\
\bottomrule
\end{tabular}
\end{table}

\noindent\textbf{Key points:} Adam includes bias correction via $\hat{\vect{m}}_t = \vect{m}_t/(1-\beta_1^t)$ and $\hat{\vect{v}}_t = \vect{v}_t/(1-\beta_2^t)$. RMSprop and Adam adapt per-parameter learning rates; momentum accelerates convergence in consistent gradient directions.

%------------------------------------------------------------------------------
\section{Step 1: Create the Neural Network}
\label{sec:capstone-network}

We'll implement a simple MLP from scratch using NumPy. This gives full control and transparency.

\begin{pythoncode}
import numpy as np

class MLP:
    """
    Multi-Layer Perceptron for MNIST classification.
    Architecture: 784 -> 256 -> 128 -> 10
    """

    def __init__(self, seed=42):
        """Initialize network with He initialization."""
        np.random.seed(seed)

        # He initialization: scale by sqrt(2/fan_in) for ReLU
        self.W1 = np.random.randn(784, 256) * np.sqrt(2.0 / 784)
        self.b1 = np.zeros((1, 256))

        self.W2 = np.random.randn(256, 128) * np.sqrt(2.0 / 256)
        self.b2 = np.zeros((1, 128))

        self.W3 = np.random.randn(128, 10) * np.sqrt(2.0 / 128)
        self.b3 = np.zeros((1, 10))

        # Cache for backward pass
        self.cache = {}

    def relu(self, x):
        """ReLU activation function."""
        return np.maximum(0, x)

    def relu_derivative(self, x):
        """Derivative of ReLU."""
        return (x > 0).astype(float)

    def softmax(self, x):
        """Numerically stable softmax."""
        exp_x = np.exp(x - np.max(x, axis=1, keepdims=True))
        return exp_x / np.sum(exp_x, axis=1, keepdims=True)

    def forward(self, X):
        """
        Forward pass through the network.

        Args:
            X: Input batch of shape (batch_size, 784)

        Returns:
            Softmax probabilities of shape (batch_size, 10)
        """
        # Layer 1
        self.cache['z1'] = X @ self.W1 + self.b1
        self.cache['a1'] = self.relu(self.cache['z1'])

        # Layer 2
        self.cache['z2'] = self.cache['a1'] @ self.W2 + self.b2
        self.cache['a2'] = self.relu(self.cache['z2'])

        # Output layer
        self.cache['z3'] = self.cache['a2'] @ self.W3 + self.b3
        self.cache['a3'] = self.softmax(self.cache['z3'])

        self.cache['X'] = X  # Save input for backward pass

        return self.cache['a3']

    def backward(self, y_true):
        """
        Backward pass using cached activations.

        Args:
            y_true: One-hot encoded labels of shape (batch_size, 10)

        Returns:
            Dictionary of gradients for all parameters
        """
        batch_size = y_true.shape[0]
        grads = {}

        # Output layer gradient (softmax + cross-entropy simplifies)
        dz3 = self.cache['a3'] - y_true  # Shape: (batch_size, 10)
        grads['W3'] = self.cache['a2'].T @ dz3 / batch_size
        grads['b3'] = np.mean(dz3, axis=0, keepdims=True)

        # Layer 2 gradients
        da2 = dz3 @ self.W3.T
        dz2 = da2 * self.relu_derivative(self.cache['z2'])
        grads['W2'] = self.cache['a1'].T @ dz2 / batch_size
        grads['b2'] = np.mean(dz2, axis=0, keepdims=True)

        # Layer 1 gradients
        da1 = dz2 @ self.W2.T
        dz1 = da1 * self.relu_derivative(self.cache['z1'])
        grads['W1'] = self.cache['X'].T @ dz1 / batch_size
        grads['b1'] = np.mean(dz1, axis=0, keepdims=True)

        return grads

    def cross_entropy_loss(self, y_pred, y_true):
        """
        Compute cross-entropy loss.

        Args:
            y_pred: Predicted probabilities (batch_size, 10)
            y_true: One-hot encoded labels (batch_size, 10)

        Returns:
            Scalar loss value
        """
        # Clip predictions to avoid log(0)
        y_pred = np.clip(y_pred, 1e-15, 1 - 1e-15)
        return -np.mean(np.sum(y_true * np.log(y_pred), axis=1))

    def get_params(self):
        """Return a list of all parameter arrays."""
        return [self.W1, self.b1, self.W2, self.b2, self.W3, self.b3]

    def set_params(self, params):
        """Set all parameters from a list."""
        self.W1, self.b1, self.W2, self.b2, self.W3, self.b3 = params

    def copy_params(self):
        """Return deep copy of parameters."""
        return [p.copy() for p in self.get_params()]
\end{pythoncode}

\begin{keyconcept}
\textbf{Implementation Details to Note:}
\begin{itemize}
    \item \textbf{He initialization}: We scale weights by $\sqrt{2/\text{fan\_in}}$ for ReLU activations. This prevents vanishing/exploding gradients at initialization.
    \item \textbf{Numerically stable softmax}: Subtracting $\max(x)$ prevents overflow when computing $e^x$.
    \item \textbf{Gradient caching}: We store intermediate activations during forward pass to use in backward pass.
    \item \textbf{Gradient averaging}: We divide by batch size to get the mean gradient.
\end{itemize}
\end{keyconcept}

%------------------------------------------------------------------------------
\section{Step 2: Implement Base Optimizer}
\label{sec:capstone-base-optimizer}

We'll create a base optimizer class that defines the interface all optimizers must follow.

\begin{pythoncode}
class Optimizer:
    """
    Base optimizer class defining the interface.

    All optimizers must implement:
    - step(model, grads): Update model parameters using gradients
    - reset(): Reset optimizer state (for fair comparisons)
    """

    def __init__(self, learning_rate):
        self.lr = learning_rate

    def step(self, model, grads):
        """
        Update model parameters using computed gradients.

        Args:
            model: The neural network with get_params()/set_params()
            grads: Dictionary of gradients {'W1': ..., 'b1': ..., etc.}
        """
        raise NotImplementedError("Subclasses must implement step()")

    def reset(self):
        """Reset optimizer state (momentum, RMSprop cache, etc.)."""
        pass  # Default: no state to reset

    def get_name(self):
        """Return optimizer name for logging."""
        return self.__class__.__name__
\end{pythoncode}

%------------------------------------------------------------------------------
\section{Step 3: Implement SGD}
\label{sec:capstone-sgd}

Vanilla SGD is the simplest optimizer---just move in the opposite direction of the gradient.

\begin{pythoncode}
class SGD(Optimizer):
    """
    Vanilla Stochastic Gradient Descent.

    Update rule: theta = theta - lr * gradient
    """

    def __init__(self, learning_rate=0.01):
        super().__init__(learning_rate)

    def step(self, model, grads):
        """Apply SGD update to all parameters."""
        param_names = ['W1', 'b1', 'W2', 'b2', 'W3', 'b3']
        params = model.get_params()

        for i, name in enumerate(param_names):
            params[i] = params[i] - self.lr * grads[name]

        model.set_params(params)
\end{pythoncode}

\begin{example}{SGD in Action}{sgd-example}
Consider a single weight $w = 0.5$ with gradient $g = 0.1$ and learning rate $\eta = 0.1$.

After one SGD step:
\[
    w_{\text{new}} = w - \eta \cdot g = 0.5 - 0.1 \times 0.1 = 0.49
\]

The weight moved slightly in the direction that decreases the loss.
\end{example}

%------------------------------------------------------------------------------
\section{Step 4: Implement SGD with Momentum}
\label{sec:capstone-momentum}

Momentum maintains a velocity vector for each parameter that accumulates gradient history.

\begin{pythoncode}
class SGDMomentum(Optimizer):
    """
    SGD with Momentum - pedagogical implementation.

    Update rule:
        v = beta * v + gradient
        theta = theta - lr * v

    NOTE: The velocity v must persist across iterations.
    Production code should use torch.optim.SGD(momentum=0.9)
    which handles this automatically via optimizer.state.
    """

    def __init__(self, learning_rate=0.01, momentum=0.9):
        super().__init__(learning_rate)
        self.beta = momentum
        self.velocity = None  # Will be initialized on first step

    def step(self, model, grads):
        """Apply momentum update to all parameters."""
        param_names = ['W1', 'b1', 'W2', 'b2', 'W3', 'b3']
        params = model.get_params()

        # Initialize velocity on first step
        if self.velocity is None:
            self.velocity = {}
            for name in param_names:
                self.velocity[name] = np.zeros_like(grads[name])

        for i, name in enumerate(param_names):
            # Update velocity: v = beta * v + gradient
            self.velocity[name] = (self.beta * self.velocity[name] +
                                   grads[name])

            # Update parameter: theta = theta - lr * v
            params[i] = params[i] - self.lr * self.velocity[name]

        model.set_params(params)

    def reset(self):
        """Reset velocity to zero."""
        self.velocity = None

    def get_name(self):
        return f"SGD+Momentum(beta={self.beta})"
\end{pythoncode}

\begin{intuition}
\textbf{Why $\beta = 0.9$?}

With $\beta = 0.9$, the effective window of gradient history is about $1/(1-\beta) = 10$ steps. Gradients from more than 10 steps ago contribute less than $0.9^{10} \approx 35\%$ of their original weight.

\begin{itemize}
    \item \textbf{Too small} ($\beta = 0.5$): Not enough smoothing, behaves almost like SGD
    \item \textbf{Too large} ($\beta = 0.99$): Too much inertia, may overshoot and oscillate
    \item \textbf{Just right} ($\beta = 0.9$): Good balance of acceleration and stability
\end{itemize}
\end{intuition}

%------------------------------------------------------------------------------
\section{Step 5: Implement RMSprop}
\label{sec:capstone-rmsprop}

RMSprop adapts learning rates per-parameter by tracking squared gradient magnitudes.

\begin{pythoncode}
class RMSprop(Optimizer):
    """
    RMSprop: Root Mean Square Propagation - pedagogical implementation.

    Update rule:
        s = rho * s + (1 - rho) * gradient^2
        theta = theta - lr * gradient / sqrt(s + epsilon)

    NOTE: The cache s must persist across iterations.
    Production code should use torch.optim.RMSprop which handles
    this automatically via optimizer.state.
    """

    def __init__(self, learning_rate=0.001, rho=0.99, epsilon=1e-8):
        super().__init__(learning_rate)
        self.rho = rho
        self.epsilon = epsilon
        self.cache = None  # Running average of squared gradients

    def step(self, model, grads):
        """Apply RMSprop update to all parameters."""
        param_names = ['W1', 'b1', 'W2', 'b2', 'W3', 'b3']
        params = model.get_params()

        # Initialize cache on first step
        if self.cache is None:
            self.cache = {}
            for name in param_names:
                self.cache[name] = np.zeros_like(grads[name])

        for i, name in enumerate(param_names):
            # Update running average of squared gradients
            self.cache[name] = (self.rho * self.cache[name] +
                               (1 - self.rho) * grads[name]**2)

            # Compute adaptive learning rate and update
            params[i] = params[i] - (self.lr * grads[name] /
                                     (np.sqrt(self.cache[name]) + self.epsilon))

        model.set_params(params)

    def reset(self):
        """Reset squared gradient cache."""
        self.cache = None

    def get_name(self):
        return f"RMSprop(rho={self.rho})"
\end{pythoncode}

\begin{example}{RMSprop Adapts to Gradient Scale}{rmsprop-example}
Consider two parameters with very different gradient magnitudes:
\begin{itemize}
    \item Parameter A: gradient $g_A = 10.0$ (steep direction)
    \item Parameter B: gradient $g_B = 0.01$ (flat direction)
\end{itemize}

After a few steps with $\rho = 0.99$, the caches become approximately $s_A \approx 100$ and $s_B \approx 0.0001$.

Effective updates (with $\eta = 0.001$):
\begin{align*}
    \Delta \theta_A &= \frac{0.001 \times 10.0}{\sqrt{100}} = 0.001 \\
    \Delta \theta_B &= \frac{0.001 \times 0.01}{\sqrt{0.0001}} = 0.001
\end{align*}

Despite 1000x different gradients, both parameters make equal-sized steps! RMSprop ``equalizes'' the optimization landscape.
\end{example}

%------------------------------------------------------------------------------
\section{Step 6: Implement Adam}
\label{sec:capstone-adam}

Adam combines the best of momentum and RMSprop, plus bias correction.

\begin{pythoncode}
class Adam(Optimizer):
    """
    Adam: Adaptive Moment Estimation - pedagogical implementation.

    Combines momentum (first moment) with RMSprop (second moment)
    and adds bias correction for early steps.

    Update rule:
        m = beta1 * m + (1 - beta1) * gradient        (momentum)
        v = beta2 * v + (1 - beta2) * gradient^2      (RMSprop)
        m_hat = m / (1 - beta1^t)                     (bias correction)
        v_hat = v / (1 - beta2^t)                     (bias correction)
        theta = theta - lr * m_hat / (sqrt(v_hat) + epsilon)

    NOTE: The moments m, v and timestep t must persist across iterations.
    Production code should use torch.optim.Adam which handles
    this automatically via optimizer.state.
    """

    def __init__(self, learning_rate=0.001, beta1=0.9, beta2=0.999,
                 epsilon=1e-8):
        super().__init__(learning_rate)
        self.beta1 = beta1
        self.beta2 = beta2
        self.epsilon = epsilon
        self.m = None  # First moment (momentum)
        self.v = None  # Second moment (RMSprop-like)
        self.t = 0     # Timestep for bias correction

    def step(self, model, grads):
        """Apply Adam update to all parameters."""
        param_names = ['W1', 'b1', 'W2', 'b2', 'W3', 'b3']
        params = model.get_params()

        # Initialize moment estimates on first step
        if self.m is None:
            self.m = {}
            self.v = {}
            for name in param_names:
                self.m[name] = np.zeros_like(grads[name])
                self.v[name] = np.zeros_like(grads[name])

        # Increment timestep
        self.t += 1

        for i, name in enumerate(param_names):
            # Update biased first moment estimate (momentum)
            self.m[name] = (self.beta1 * self.m[name] +
                          (1 - self.beta1) * grads[name])

            # Update biased second moment estimate (RMSprop)
            self.v[name] = (self.beta2 * self.v[name] +
                          (1 - self.beta2) * grads[name]**2)

            # Compute bias-corrected estimates
            m_hat = self.m[name] / (1 - self.beta1**self.t)
            v_hat = self.v[name] / (1 - self.beta2**self.t)

            # Update parameters
            params[i] = params[i] - (self.lr * m_hat /
                                     (np.sqrt(v_hat) + self.epsilon))

        model.set_params(params)

    def reset(self):
        """Reset moment estimates and timestep."""
        self.m = None
        self.v = None
        self.t = 0

    def get_name(self):
        return f"Adam(beta1={self.beta1}, beta2={self.beta2})"
\end{pythoncode}

\begin{important}
\textbf{The Bias Correction Factor in Detail}

At $t=1$ with default $\beta_1 = 0.9$ and $\beta_2 = 0.999$:
\begin{align*}
    \text{Correction for } m: \quad \frac{1}{1 - 0.9^1} &= \frac{1}{0.1} = 10\times \\
    \text{Correction for } v: \quad \frac{1}{1 - 0.999^1} &= \frac{1}{0.001} = 1000\times
\end{align*}

Without this correction, the first few updates would be far too small because the moment estimates start at zero. By $t=100$:
\begin{align*}
    \text{Correction for } m: \quad \frac{1}{1 - 0.9^{100}} &\approx 1.0000 \\
    \text{Correction for } v: \quad \frac{1}{1 - 0.999^{100}} &\approx 1.105
\end{align*}

The correction quickly becomes negligible.
\end{important}

%------------------------------------------------------------------------------
\section{Step 7: Training Loop}
\label{sec:capstone-training}

Now we need a training function that works with any optimizer and tracks metrics.

\begin{pythoncode}
def load_mnist():
    """
    Load MNIST dataset using sklearn.
    Returns train/val/test splits with normalized pixel values.

    Split ratios: 80% train, 10% validation, 10% test
    - Validation: used for optimizer comparison during training
    - Test: used only once at the end for final evaluation
    """
    from sklearn.datasets import fetch_openml
    from sklearn.model_selection import train_test_split

    # Load MNIST
    print("Loading MNIST dataset...")
    mnist = fetch_openml('mnist_784', version=1, as_frame=False)
    X, y = mnist.data, mnist.target.astype(int)

    # Normalize to [0, 1]
    X = X / 255.0

    # One-hot encode labels
    y_onehot = np.zeros((len(y), 10))
    y_onehot[np.arange(len(y)), y] = 1

    # First split: separate test set (10%)
    X_temp, X_test, y_temp, y_test = train_test_split(
        X, y_onehot, test_size=0.1, random_state=42, stratify=y
    )

    # Second split: separate validation from training (10% of original = 11.1% of temp)
    X_train, X_val, y_train, y_val = train_test_split(
        X_temp, y_temp, test_size=0.1111, random_state=42,
        stratify=np.argmax(y_temp, axis=1)
    )

    print(f"Training samples: {len(X_train)}")
    print(f"Validation samples: {len(X_val)}")
    print(f"Test samples: {len(X_test)}")

    return X_train, X_val, X_test, y_train, y_val, y_test


def compute_accuracy(model, X, y):
    """
    Compute classification accuracy.

    Args:
        model: The trained MLP
        X: Input features
        y: One-hot encoded labels

    Returns:
        Accuracy as a float in [0, 1]
    """
    y_pred = model.forward(X)
    pred_classes = np.argmax(y_pred, axis=1)
    true_classes = np.argmax(y, axis=1)
    return np.mean(pred_classes == true_classes)


def create_batches(X, y, batch_size, shuffle=True):
    """
    Create mini-batches from data.

    Args:
        X: Input features
        y: Labels
        batch_size: Size of each batch
        shuffle: Whether to shuffle data

    Yields:
        (X_batch, y_batch) tuples
    """
    n_samples = len(X)
    indices = np.arange(n_samples)

    if shuffle:
        np.random.shuffle(indices)

    for start in range(0, n_samples, batch_size):
        end = min(start + batch_size, n_samples)
        batch_idx = indices[start:end]
        yield X[batch_idx], y[batch_idx]


def train(model, optimizer, X_train, y_train, X_val, y_val,
          epochs=20, batch_size=128, verbose=True):
    """
    Train the model using the given optimizer.

    Args:
        model: The MLP to train
        optimizer: The optimizer to use
        X_train, y_train: Training data
        X_val, y_val: Validation data (for monitoring, NOT test data)
        epochs: Number of training epochs
        batch_size: Mini-batch size
        verbose: Whether to print progress

    Returns:
        Dictionary with training history

    Note: We evaluate on validation set each epoch for optimizer comparison.
          The test set is held out and used only for final evaluation.
    """
    history = {
        'train_loss': [],
        'train_acc': [],
        'val_loss': [],
        'val_acc': [],
        'epoch_times': []
    }

    import time

    for epoch in range(epochs):
        epoch_start = time.time()
        epoch_losses = []

        # Training loop
        for X_batch, y_batch in create_batches(X_train, y_train, batch_size):
            # Forward pass
            y_pred = model.forward(X_batch)
            loss = model.cross_entropy_loss(y_pred, y_batch)
            epoch_losses.append(loss)

            # Backward pass
            grads = model.backward(y_batch)

            # Optimizer step
            optimizer.step(model, grads)

        epoch_time = time.time() - epoch_start

        # Compute metrics
        train_loss = np.mean(epoch_losses)
        train_acc = compute_accuracy(model, X_train, y_train)

        # Validation metrics (NOT test metrics!)
        val_pred = model.forward(X_val)
        val_loss = model.cross_entropy_loss(val_pred, y_val)
        val_acc = compute_accuracy(model, X_val, y_val)

        # Record history
        history['train_loss'].append(train_loss)
        history['train_acc'].append(train_acc)
        history['val_loss'].append(val_loss)
        history['val_acc'].append(val_acc)
        history['epoch_times'].append(epoch_time)

        if verbose:
            print(f"Epoch {epoch+1:3d}/{epochs} | "
                  f"Loss: {train_loss:.4f} | "
                  f"Train Acc: {train_acc:.4f} | "
                  f"Val Acc: {val_acc:.4f} | "
                  f"Time: {epoch_time:.2f}s")

    return history
\end{pythoncode}

%------------------------------------------------------------------------------
\section{Step 8: Run Experiments}
\label{sec:capstone-experiments}

Now let's run our showdown! We'll train with each optimizer using learning rates typical for each method.

\begin{pythoncode}
def run_optimizer_showdown():
    """
    Main experiment: compare all optimizers on MNIST.

    Uses proper train/val/test split:
    - Training: 80% for parameter updates
    - Validation: 10% for optimizer comparison during training
    - Test: 10% for final evaluation (used only once at the end)
    """
    # Load data with proper three-way split
    X_train, X_val, X_test, y_train, y_val, y_test = load_mnist()

    # Define optimizers with typical learning rates
    optimizers = [
        SGD(learning_rate=0.1),
        SGDMomentum(learning_rate=0.01, momentum=0.9),
        RMSprop(learning_rate=0.001, rho=0.99),
        Adam(learning_rate=0.001, beta1=0.9, beta2=0.999),
    ]

    # Store initial parameters for fair comparison
    # (all optimizers start from the same initialization)
    base_model = MLP(seed=42)
    initial_params = base_model.copy_params()

    # Run experiments
    all_results = {}
    trained_models = {}  # Store models for final test evaluation

    for opt in optimizers:
        print(f"\n{'='*60}")
        print(f"Training with {opt.get_name()}")
        print(f"{'='*60}")

        # Reset model to initial parameters
        model = MLP(seed=42)
        model.set_params([p.copy() for p in initial_params])

        # Reset optimizer state
        opt.reset()

        # Train using validation set for monitoring (NOT test set!)
        history = train(
            model, opt,
            X_train, y_train,
            X_val, y_val,  # Validation set, not test set
            epochs=20,
            batch_size=128
        )

        all_results[opt.get_name()] = history
        trained_models[opt.get_name()] = model

    # Final evaluation on held-out test set (done only once!)
    print(f"\n{'='*60}")
    print("FINAL TEST SET EVALUATION (single evaluation)")
    print(f"{'='*60}")
    for name, model in trained_models.items():
        test_acc = compute_accuracy(model, X_test, y_test)
        all_results[name]['final_test_acc'] = test_acc
        print(f"{name.split('(')[0]:<20} Test Accuracy: {test_acc*100:.2f}%")

    return all_results


# Run the experiments
if __name__ == "__main__":
    results = run_optimizer_showdown()
\end{pythoncode}

\begin{keyconcept}
\textbf{Why Different Learning Rates?}

Each optimizer operates best at different scales:
\begin{itemize}
    \item \textbf{SGD} ($\eta = 0.1$): Needs larger learning rate because it has no momentum or adaptation
    \item \textbf{SGD+Momentum} ($\eta = 0.01$): Smaller rate because momentum amplifies steps
    \item \textbf{RMSprop} ($\eta = 0.001$): Small rate because adaptive scaling can make large steps
    \item \textbf{Adam} ($\eta = 0.001$): Same as RMSprop; this is the well-known default
\end{itemize}

These are ``tuned'' starting points. In practice, you would search around these values.
\end{keyconcept}

%------------------------------------------------------------------------------
\section{Step 9: Analyze Results}
\label{sec:capstone-analysis}

Let's create visualizations to compare the optimizers.

\begin{pythoncode}
import matplotlib.pyplot as plt

def plot_results(all_results):
    """
    Create comparison plots for all optimizers.

    Note: Plots show validation accuracy (used during training for comparison).
    Final test accuracy is reported separately and evaluated only once.
    """
    fig, axes = plt.subplots(2, 2, figsize=(14, 10))

    # Color scheme
    colors = {
        'SGD': '#1f77b4',
        'SGD+Momentum': '#ff7f0e',
        'RMSprop': '#2ca02c',
        'Adam': '#d62728'
    }

    # Helper to get color (check more specific names first)
    def get_color(name):
        # Check 'SGD+Momentum' before 'SGD' to avoid substring mismatch
        if 'Momentum' in name:
            return colors['SGD+Momentum']
        for key in ['Adam', 'RMSprop', 'SGD']:
            if key in name:
                return colors[key]
        return 'black'

    epochs = range(1, 21)

    # Plot 1: Training Loss
    ax1 = axes[0, 0]
    for name, history in all_results.items():
        ax1.plot(epochs, history['train_loss'],
                label=name.split('(')[0],  # Shorter name
                color=get_color(name), linewidth=2)
    ax1.set_xlabel('Epoch')
    ax1.set_ylabel('Training Loss')
    ax1.set_title('Training Loss Comparison')
    ax1.legend()
    ax1.grid(True, alpha=0.3)
    ax1.set_yscale('log')

    # Plot 2: Validation Accuracy (NOT test accuracy!)
    ax2 = axes[0, 1]
    for name, history in all_results.items():
        ax2.plot(epochs, [a * 100 for a in history['val_acc']],
                label=name.split('(')[0],
                color=get_color(name), linewidth=2)
    ax2.set_xlabel('Epoch')
    ax2.set_ylabel('Validation Accuracy (%)')
    ax2.set_title('Validation Accuracy Comparison')
    ax2.legend()
    ax2.grid(True, alpha=0.3)
    ax2.axhline(y=95, color='gray', linestyle='--', alpha=0.5)

    # Plot 3: Training vs Validation Accuracy (generalization gap)
    ax3 = axes[1, 0]
    for name, history in all_results.items():
        train_acc = [a * 100 for a in history['train_acc']]
        val_acc = [a * 100 for a in history['val_acc']]
        gap = [t - v for t, v in zip(train_acc, val_acc)]
        ax3.plot(epochs, gap,
                label=name.split('(')[0],
                color=get_color(name), linewidth=2)
    ax3.set_xlabel('Epoch')
    ax3.set_ylabel('Train - Validation Accuracy (%)')
    ax3.set_title('Generalization Gap')
    ax3.legend()
    ax3.grid(True, alpha=0.3)

    # Plot 4: Convergence speed (epochs to reach 95% validation accuracy)
    ax4 = axes[1, 1]
    final_metrics = []
    for name, history in all_results.items():
        val_acc = history['val_acc']
        # Find first epoch to reach 95% validation accuracy
        epochs_to_95 = None
        for i, acc in enumerate(val_acc):
            if acc >= 0.95:
                epochs_to_95 = i + 1
                break

        final_metrics.append({
            'name': name.split('(')[0],
            'final_val_acc': val_acc[-1] * 100,
            'epochs_to_95': epochs_to_95 if epochs_to_95 else '>20',
            'total_time': sum(history['epoch_times'])
        })

    # Create bar chart showing epochs to 95% (convergence speed)
    names = [m['name'] for m in final_metrics]
    epochs_values = [m['epochs_to_95'] if isinstance(m['epochs_to_95'], int)
                     else 21 for m in final_metrics]
    bar_colors = [get_color(m['name']) for m in final_metrics]

    bars = ax4.bar(names, epochs_values, color=bar_colors)
    ax4.set_ylabel('Epochs to 95% Validation Accuracy')
    ax4.set_title('Convergence Speed (fewer epochs = faster)')
    ax4.set_ylim(0, 25)

    # Add value labels on bars
    for bar, m in zip(bars, final_metrics):
        label = str(m['epochs_to_95'])
        ax4.text(bar.get_x() + bar.get_width()/2, bar.get_height() + 0.3,
                label, ha='center', va='bottom', fontsize=10)

    plt.tight_layout()
    plt.savefig('optimizer_comparison.png', dpi=150, bbox_inches='tight')
    plt.show()

    return final_metrics


def print_summary_table(final_metrics, all_results):
    """
    Print a summary table of results.

    Shows both validation accuracy (used during training) and
    final test accuracy (evaluated only once at the end).
    """
    print("\n" + "="*80)
    print("OPTIMIZER SHOWDOWN RESULTS")
    print("="*80)

    print(f"\n{'Optimizer':<20} {'Val Acc':<12} {'Test Acc':<12} "
          f"{'Epochs to 95%':<15} {'Total Time':<12}")
    print("-"*80)

    for name, history in all_results.items():
        val_acc = history['val_acc']
        final_val_acc = val_acc[-1] * 100
        # Final test accuracy (evaluated only once)
        final_test_acc = history.get('final_test_acc', 0) * 100
        total_time = sum(history['epoch_times'])

        # Find epochs to 95% validation accuracy
        epochs_to_95 = "N/A"
        for i, acc in enumerate(val_acc):
            if acc >= 0.95:
                epochs_to_95 = str(i + 1)
                break

        short_name = name.split('(')[0]
        print(f"{short_name:<20} {final_val_acc:<12.2f}% {final_test_acc:<12.2f}% "
              f"{epochs_to_95:<15} {total_time:<12.2f}s")

    print("\nNote: Val Acc = validation accuracy (monitored during training)")
    print("      Test Acc = held-out test accuracy (evaluated only once)")


# After running experiments:
# final_metrics = plot_results(results)
# print_summary_table(final_metrics, results)
\end{pythoncode}

\subsection{Learning Rate Sensitivity Analysis}

To understand how each optimizer handles different learning rates, let's sweep across values.

\begin{pythoncode}
def learning_rate_sensitivity(X_train, y_train, X_val, y_val):
    """
    Test each optimizer across a range of learning rates.

    Uses validation set for comparison (not test set).
    """
    # Learning rate multipliers (relative to default)
    multipliers = [0.1, 0.3, 1.0, 3.0, 10.0]

    # Base learning rates
    base_lrs = {
        'SGD': 0.1,
        'Momentum': 0.01,
        'RMSprop': 0.001,
        'Adam': 0.001
    }

    sensitivity_results = {name: {'lrs': [], 'accs': []}
                          for name in base_lrs}

    base_model = MLP(seed=42)
    initial_params = base_model.copy_params()

    for mult in multipliers:
        print(f"\nTesting learning rate multiplier: {mult}x")

        for name, base_lr in base_lrs.items():
            lr = base_lr * mult

            # Create optimizer
            if name == 'SGD':
                opt = SGD(learning_rate=lr)
            elif name == 'Momentum':
                opt = SGDMomentum(learning_rate=lr, momentum=0.9)
            elif name == 'RMSprop':
                opt = RMSprop(learning_rate=lr)
            else:  # Adam
                opt = Adam(learning_rate=lr)

            # Train for 10 epochs
            model = MLP(seed=42)
            model.set_params([p.copy() for p in initial_params])
            opt.reset()

            history = train(model, opt, X_train, y_train,
                          X_val, y_val, epochs=10,
                          batch_size=128, verbose=False)

            final_acc = history['val_acc'][-1]
            sensitivity_results[name]['lrs'].append(lr)
            sensitivity_results[name]['accs'].append(final_acc)

            print(f"  {name}: lr={lr:.4f}, val_acc={final_acc:.4f}")

    return sensitivity_results


def plot_sensitivity(sensitivity_results):
    """
    Plot learning rate sensitivity curves.
    """
    plt.figure(figsize=(10, 6))

    colors = {
        'SGD': '#1f77b4',
        'Momentum': '#ff7f0e',
        'RMSprop': '#2ca02c',
        'Adam': '#d62728'
    }

    for name, data in sensitivity_results.items():
        plt.plot(data['lrs'], [a * 100 for a in data['accs']],
                'o-', label=name, color=colors[name],
                linewidth=2, markersize=8)

    plt.xlabel('Learning Rate')
    plt.ylabel('Validation Accuracy after 10 epochs (%)')
    plt.title('Learning Rate Sensitivity')
    plt.xscale('log')
    plt.legend()
    plt.grid(True, alpha=0.3)
    plt.savefig('lr_sensitivity.png', dpi=150, bbox_inches='tight')
    plt.show()
\end{pythoncode}

%------------------------------------------------------------------------------
\section{Key Findings}
\label{sec:capstone-findings}

After running the experiments, here's what you should observe (based on \textit{validation} accuracy, not test accuracy---the test set is evaluated only once at the end):

\begin{keyconcept}
\textbf{Convergence Speed:}
\begin{itemize}
    \item \textbf{Adam} converges fastest, typically reaching 95\% validation accuracy in 3--5 epochs
    \item \textbf{RMSprop} is close behind, often 4--6 epochs
    \item \textbf{SGD+Momentum} is moderate, usually 6--10 epochs
    \item \textbf{Vanilla SGD} is slowest, often needing 10+ epochs or failing to reach 95\%
\end{itemize}
\end{keyconcept}

\begin{keyconcept}
\textbf{Final Accuracy:}
All well-tuned optimizers achieve similar final validation accuracy (97--98\%) on MNIST. The differences are small---this is a ``solved'' problem. On harder datasets, differences become more pronounced. Final test accuracy (evaluated once) should be similar to validation accuracy if the methodology is correct.
\end{keyconcept}

\begin{keyconcept}
\textbf{Training Stability:}
\begin{itemize}
    \item \textbf{Adam} shows the smoothest loss curves
    \item \textbf{RMSprop} is almost as smooth
    \item \textbf{SGD+Momentum} shows some oscillation, especially early
    \item \textbf{Vanilla SGD} can be quite noisy, especially with larger learning rates
\end{itemize}
\end{keyconcept}

\begin{keyconcept}
\textbf{Learning Rate Sensitivity:}
\begin{itemize}
    \item \textbf{Adam} is most robust---works well across a wide range
    \item \textbf{RMSprop} is also fairly robust
    \item \textbf{SGD+Momentum} is sensitive---too high diverges, too low stalls
    \item \textbf{Vanilla SGD} is most sensitive---narrow range of working values
\end{itemize}
\end{keyconcept}

\begin{table}[htbp]
\centering
\caption{Summary of optimizer characteristics on MNIST (validation set metrics)}
\label{tab:optimizer-summary}
\begin{tabular}{@{}lcccc@{}}
\toprule
\textbf{Metric} & \textbf{SGD} & \textbf{Momentum} & \textbf{RMSprop} & \textbf{Adam} \\
\midrule
Convergence Speed & Slow & Medium & Fast & Fastest \\
Final Val Accuracy & 97.0\% & 97.8\% & 97.9\% & 98.0\% \\
Stability & Low & Medium & High & Highest \\
LR Sensitivity & High & Medium & Low & Lowest \\
Memory Overhead & None & 1x params & 1x params & 2x params \\
\bottomrule
\end{tabular}
\end{table}

%------------------------------------------------------------------------------
\section{When to Use Each Optimizer}
\label{sec:capstone-recommendations}

Based on our experiments and broader experience, here are practical recommendations:

\begin{principle}
\textbf{Adam/AdamW: The Default Choice}
\begin{itemize}
    \item Use for \textbf{quick prototyping}---works well out of the box
    \item Use for \textbf{Transformers and NLP}---standard choice, well-tested
    \item Use for \textbf{training from scratch} when you don't want to tune learning rates
    \item \textbf{Typical settings}: lr=0.001, betas=(0.9, 0.999), weight\_decay=0.01
\end{itemize}
\end{principle}

\begin{principle}
\textbf{SGD+Momentum: For Maximum Performance}
\begin{itemize}
    \item Use for \textbf{CNNs on image classification}---often generalizes better
    \item Use when \textbf{squeezing out last few percent accuracy matters}
    \item Requires \textbf{more careful tuning} of learning rate and schedule
    \item \textbf{Typical settings}: lr=0.1 with cosine/step decay, momentum=0.9
\end{itemize}
\end{principle}

\begin{principle}
\textbf{RMSprop: For Recurrent Networks}
\begin{itemize}
    \item Use for \textbf{RNNs and LSTMs}---handles exploding gradients well
    \item Use when Adam behaves unexpectedly
    \item \textbf{Typical settings}: lr=0.001, rho=0.9 or 0.99
\end{itemize}
\end{principle}

\begin{principle}
\textbf{Vanilla SGD: For Research and Understanding}
\begin{itemize}
    \item Use when \textbf{simplicity matters} for reproducibility
    \item Use for \textbf{theoretical analysis}---easier to reason about
    \item Generally \textbf{not recommended} for practical training
\end{itemize}
\end{principle}

\begin{perspective}
\textbf{The Generalization Debate}

There's ongoing research about whether SGD finds ``better'' minima than Adam. The hypothesis is that SGD's noise helps it find wider minima that generalize better. In practice:
\begin{itemize}
    \item On \textbf{ImageNet with ResNets}, SGD+Momentum often beats Adam by 0.5--1\%
    \item On \textbf{Transformers}, AdamW consistently wins
    \item On \textbf{MNIST and CIFAR}, differences are minimal with proper tuning
\end{itemize}
The ``best'' optimizer depends on your specific problem, model, and how much time you have for tuning.
\end{perspective}

%------------------------------------------------------------------------------
\section{Extensions}
\label{sec:capstone-extensions}

Ready to go further? Here are challenging extensions:

\subsection{Extension 1: Learning Rate Scheduling}

Add learning rate schedules to your training loop:

\begin{pythoncode}
def cosine_schedule(epoch, total_epochs, initial_lr, min_lr=1e-6):
    """Cosine annealing learning rate schedule."""
    return min_lr + 0.5 * (initial_lr - min_lr) * \
           (1 + np.cos(np.pi * epoch / total_epochs))

def step_decay_schedule(epoch, initial_lr, drop_factor=0.1,
                        drop_epochs=[30, 60, 80]):
    """Step decay learning rate schedule."""
    lr = initial_lr
    for drop_epoch in drop_epochs:
        if epoch >= drop_epoch:
            lr *= drop_factor
    return lr

def warmup_schedule(epoch, warmup_epochs, initial_lr):
    """Linear warmup for the first few epochs."""
    if epoch < warmup_epochs:
        return initial_lr * (epoch + 1) / warmup_epochs
    return initial_lr
\end{pythoncode}

\textbf{Challenge}: Implement training with cosine annealing and compare results. Does scheduling help Adam as much as SGD?

\subsection{Extension 2: AdamW Implementation}

Implement AdamW with decoupled weight decay:

\begin{pythoncode}
class AdamW(Adam):
    """
    AdamW: Adam with decoupled weight decay.

    The key difference: weight decay is applied directly to parameters,
    not through the gradient. This leads to better generalization.
    """

    def __init__(self, learning_rate=0.001, beta1=0.9, beta2=0.999,
                 epsilon=1e-8, weight_decay=0.01):
        super().__init__(learning_rate, beta1, beta2, epsilon)
        self.weight_decay = weight_decay

    def step(self, model, grads):
        """Apply AdamW update with decoupled weight decay."""
        param_names = ['W1', 'b1', 'W2', 'b2', 'W3', 'b3']
        params = model.get_params()

        if self.m is None:
            self.m = {}
            self.v = {}
            for name in param_names:
                self.m[name] = np.zeros_like(grads[name])
                self.v[name] = np.zeros_like(grads[name])

        self.t += 1

        for i, name in enumerate(param_names):
            # Standard Adam moment updates
            self.m[name] = self.beta1 * self.m[name] + \
                          (1 - self.beta1) * grads[name]
            self.v[name] = self.beta2 * self.v[name] + \
                          (1 - self.beta2) * grads[name]**2

            m_hat = self.m[name] / (1 - self.beta1**self.t)
            v_hat = self.v[name] / (1 - self.beta2**self.t)

            # Adam update
            adam_update = self.lr * m_hat / (np.sqrt(v_hat) + self.epsilon)

            # Decoupled weight decay (only for weights, not biases)
            if 'W' in name:  # Apply weight decay only to weight matrices
                weight_decay_term = self.lr * self.weight_decay * params[i]
                params[i] = params[i] - adam_update - weight_decay_term
            else:
                params[i] = params[i] - adam_update

        model.set_params(params)

    def get_name(self):
        return f"AdamW(wd={self.weight_decay})"
\end{pythoncode}

\textbf{Challenge}: Compare Adam with L2 regularization vs. AdamW. When does the difference matter most?

\subsection{Extension 3: Harder Dataset}

MNIST is relatively easy. Try CIFAR-10 for a more challenging benchmark:

\begin{pythoncode}
def load_cifar10():
    """
    Load CIFAR-10 and preprocess for our MLP.
    Note: CNNs work much better on CIFAR-10, but this tests optimizers.

    Returns train/val/test splits (80%/10%/10%) with proper shuffling.
    """
    from sklearn.datasets import fetch_openml
    from sklearn.model_selection import train_test_split

    # Note: OpenML uses 'CIFAR_10' (with underscore), not 'cifar-10'
    print("Loading CIFAR-10...")
    cifar = fetch_openml('CIFAR_10', version=1, as_frame=False)
    X, y = cifar.data, cifar.target.astype(int)

    # Normalize
    X = X / 255.0

    # One-hot encode
    y_onehot = np.zeros((len(y), 10))
    y_onehot[np.arange(len(y)), y] = 1

    # Proper stratified split (don't just take first 50000!)
    # First split: separate test set (10%)
    X_temp, X_test, y_temp, y_test = train_test_split(
        X, y_onehot, test_size=0.1, random_state=42, stratify=y
    )

    # Second split: separate validation from training
    X_train, X_val, y_train, y_val = train_test_split(
        X_temp, y_temp, test_size=0.1111, random_state=42,
        stratify=np.argmax(y_temp, axis=1)
    )

    print(f"Training samples: {len(X_train)}")
    print(f"Validation samples: {len(X_val)}")
    print(f"Test samples: {len(X_test)}")

    return X_train, X_val, X_test, y_train, y_val, y_test
\end{pythoncode}

You'll need to adjust the network (input size 3072 for 32x32x3 images) and likely add dropout or other regularization.

\subsection{Extension 4: Gradient Clipping}

Implement gradient clipping to stabilize training with large learning rates:

\begin{pythoncode}
def clip_gradients(grads, max_norm=1.0):
    """
    Clip gradients by global norm.

    If ||grads|| > max_norm, scale all gradients so ||grads|| = max_norm.
    """
    # Compute global norm
    total_norm = 0
    for name, grad in grads.items():
        total_norm += np.sum(grad ** 2)
    total_norm = np.sqrt(total_norm)

    # Clip if necessary
    if total_norm > max_norm:
        scale = max_norm / (total_norm + 1e-8)
        grads = {name: grad * scale for name, grad in grads.items()}

    return grads, total_norm
\end{pythoncode}

\textbf{Challenge}: Does gradient clipping allow SGD to work with larger learning rates?

%------------------------------------------------------------------------------
\section{Complete Solution}
\label{sec:capstone-solution}

Here is the complete, runnable code combining all components:

\begin{pythoncode}
"""
Optimizer Showdown: Complete Implementation
Comparing SGD, Momentum, RMSprop, and Adam on MNIST

Run this script to reproduce all experiments:
    python optimizer_showdown.py
"""

import numpy as np
import matplotlib.pyplot as plt
import time


# ==============================================================================
# Neural Network Implementation
# ==============================================================================

class MLP:
    """Multi-Layer Perceptron: 784 -> 256 -> 128 -> 10"""

    def __init__(self, seed=42):
        np.random.seed(seed)
        self.W1 = np.random.randn(784, 256) * np.sqrt(2.0 / 784)
        self.b1 = np.zeros((1, 256))
        self.W2 = np.random.randn(256, 128) * np.sqrt(2.0 / 256)
        self.b2 = np.zeros((1, 128))
        self.W3 = np.random.randn(128, 10) * np.sqrt(2.0 / 128)
        self.b3 = np.zeros((1, 10))
        self.cache = {}

    def relu(self, x):
        return np.maximum(0, x)

    def relu_derivative(self, x):
        return (x > 0).astype(float)

    def softmax(self, x):
        exp_x = np.exp(x - np.max(x, axis=1, keepdims=True))
        return exp_x / np.sum(exp_x, axis=1, keepdims=True)

    def forward(self, X):
        self.cache['z1'] = X @ self.W1 + self.b1
        self.cache['a1'] = self.relu(self.cache['z1'])
        self.cache['z2'] = self.cache['a1'] @ self.W2 + self.b2
        self.cache['a2'] = self.relu(self.cache['z2'])
        self.cache['z3'] = self.cache['a2'] @ self.W3 + self.b3
        self.cache['a3'] = self.softmax(self.cache['z3'])
        self.cache['X'] = X
        return self.cache['a3']

    def backward(self, y_true):
        batch_size = y_true.shape[0]
        grads = {}

        dz3 = self.cache['a3'] - y_true
        grads['W3'] = self.cache['a2'].T @ dz3 / batch_size
        grads['b3'] = np.mean(dz3, axis=0, keepdims=True)

        da2 = dz3 @ self.W3.T
        dz2 = da2 * self.relu_derivative(self.cache['z2'])
        grads['W2'] = self.cache['a1'].T @ dz2 / batch_size
        grads['b2'] = np.mean(dz2, axis=0, keepdims=True)

        da1 = dz2 @ self.W2.T
        dz1 = da1 * self.relu_derivative(self.cache['z1'])
        grads['W1'] = self.cache['X'].T @ dz1 / batch_size
        grads['b1'] = np.mean(dz1, axis=0, keepdims=True)

        return grads

    def cross_entropy_loss(self, y_pred, y_true):
        y_pred = np.clip(y_pred, 1e-15, 1 - 1e-15)
        return -np.mean(np.sum(y_true * np.log(y_pred), axis=1))

    def get_params(self):
        return [self.W1, self.b1, self.W2, self.b2, self.W3, self.b3]

    def set_params(self, params):
        self.W1, self.b1, self.W2, self.b2, self.W3, self.b3 = params

    def copy_params(self):
        return [p.copy() for p in self.get_params()]


# ==============================================================================
# Optimizer Implementations
# ==============================================================================

class Optimizer:
    def __init__(self, learning_rate):
        self.lr = learning_rate

    def step(self, model, grads):
        raise NotImplementedError

    def reset(self):
        pass

    def get_name(self):
        return self.__class__.__name__


class SGD(Optimizer):
    def step(self, model, grads):
        param_names = ['W1', 'b1', 'W2', 'b2', 'W3', 'b3']
        params = model.get_params()
        for i, name in enumerate(param_names):
            params[i] = params[i] - self.lr * grads[name]
        model.set_params(params)


class SGDMomentum(Optimizer):
    def __init__(self, learning_rate=0.01, momentum=0.9):
        super().__init__(learning_rate)
        self.beta = momentum
        self.velocity = None

    def step(self, model, grads):
        param_names = ['W1', 'b1', 'W2', 'b2', 'W3', 'b3']
        params = model.get_params()

        if self.velocity is None:
            self.velocity = {n: np.zeros_like(grads[n]) for n in param_names}

        for i, name in enumerate(param_names):
            self.velocity[name] = self.beta * self.velocity[name] + grads[name]
            params[i] = params[i] - self.lr * self.velocity[name]
        model.set_params(params)

    def reset(self):
        self.velocity = None

    def get_name(self):
        return f"SGD+Momentum(b={self.beta})"


class RMSprop(Optimizer):
    def __init__(self, learning_rate=0.001, rho=0.99, epsilon=1e-8):
        super().__init__(learning_rate)
        self.rho = rho
        self.epsilon = epsilon
        self.cache = None

    def step(self, model, grads):
        param_names = ['W1', 'b1', 'W2', 'b2', 'W3', 'b3']
        params = model.get_params()

        if self.cache is None:
            self.cache = {n: np.zeros_like(grads[n]) for n in param_names}

        for i, name in enumerate(param_names):
            self.cache[name] = (self.rho * self.cache[name] +
                               (1 - self.rho) * grads[name]**2)
            params[i] = params[i] - (self.lr * grads[name] /
                                     (np.sqrt(self.cache[name]) + self.epsilon))
        model.set_params(params)

    def reset(self):
        self.cache = None

    def get_name(self):
        return f"RMSprop(rho={self.rho})"


class Adam(Optimizer):
    def __init__(self, learning_rate=0.001, beta1=0.9, beta2=0.999,
                 epsilon=1e-8):
        super().__init__(learning_rate)
        self.beta1, self.beta2, self.epsilon = beta1, beta2, epsilon
        self.m, self.v, self.t = None, None, 0

    def step(self, model, grads):
        param_names = ['W1', 'b1', 'W2', 'b2', 'W3', 'b3']
        params = model.get_params()

        if self.m is None:
            self.m = {n: np.zeros_like(grads[n]) for n in param_names}
            self.v = {n: np.zeros_like(grads[n]) for n in param_names}

        self.t += 1

        for i, name in enumerate(param_names):
            self.m[name] = self.beta1 * self.m[name] + (1-self.beta1) * grads[name]
            self.v[name] = self.beta2 * self.v[name] + (1-self.beta2) * grads[name]**2
            m_hat = self.m[name] / (1 - self.beta1**self.t)
            v_hat = self.v[name] / (1 - self.beta2**self.t)
            params[i] = params[i] - self.lr * m_hat / (np.sqrt(v_hat) + self.epsilon)
        model.set_params(params)

    def reset(self):
        self.m, self.v, self.t = None, None, 0

    def get_name(self):
        return f"Adam(b1={self.beta1},b2={self.beta2})"


# ==============================================================================
# Training Utilities
# ==============================================================================

def load_mnist():
    """Load MNIST with proper train/val/test split (80%/10%/10%)."""
    from sklearn.datasets import fetch_openml
    from sklearn.model_selection import train_test_split

    print("Loading MNIST...")
    mnist = fetch_openml('mnist_784', version=1, as_frame=False)
    X, y = mnist.data / 255.0, mnist.target.astype(int)

    y_onehot = np.zeros((len(y), 10))
    y_onehot[np.arange(len(y)), y] = 1

    # First split: separate test set (10%)
    X_temp, X_test, y_temp, y_test = train_test_split(
        X, y_onehot, test_size=0.1, random_state=42, stratify=y)
    # Second split: separate validation from training
    X_train, X_val, y_train, y_val = train_test_split(
        X_temp, y_temp, test_size=0.1111, random_state=42,
        stratify=np.argmax(y_temp, axis=1))

    return X_train, X_val, X_test, y_train, y_val, y_test


def compute_accuracy(model, X, y):
    y_pred = model.forward(X)
    return np.mean(np.argmax(y_pred, axis=1) == np.argmax(y, axis=1))


def create_batches(X, y, batch_size):
    indices = np.random.permutation(len(X))
    for start in range(0, len(X), batch_size):
        idx = indices[start:start + batch_size]
        yield X[idx], y[idx]


def train(model, optimizer, X_train, y_train, X_val, y_val,
          epochs=20, batch_size=128, verbose=True):
    """Train using validation set for monitoring (NOT test set)."""
    history = {'train_loss': [], 'train_acc': [],
               'val_loss': [], 'val_acc': [], 'epoch_times': []}

    for epoch in range(epochs):
        t0 = time.time()
        losses = []

        for X_batch, y_batch in create_batches(X_train, y_train, batch_size):
            y_pred = model.forward(X_batch)
            losses.append(model.cross_entropy_loss(y_pred, y_batch))
            grads = model.backward(y_batch)
            optimizer.step(model, grads)

        epoch_time = time.time() - t0
        train_loss, train_acc = np.mean(losses), compute_accuracy(model, X_train, y_train)
        val_pred = model.forward(X_val)
        val_loss = model.cross_entropy_loss(val_pred, y_val)
        val_acc = compute_accuracy(model, X_val, y_val)

        history['train_loss'].append(train_loss)
        history['train_acc'].append(train_acc)
        history['val_loss'].append(val_loss)
        history['val_acc'].append(val_acc)
        history['epoch_times'].append(epoch_time)

        if verbose:
            print(f"Epoch {epoch+1:2d} | Loss: {train_loss:.4f} | "
                  f"Train: {train_acc:.4f} | Val: {val_acc:.4f}")

    return history


# ==============================================================================
# Main Experiment
# ==============================================================================

def run_showdown():
    """Compare optimizers using proper train/val/test methodology."""
    X_train, X_val, X_test, y_train, y_val, y_test = load_mnist()

    optimizers = [
        SGD(learning_rate=0.1),
        SGDMomentum(learning_rate=0.01, momentum=0.9),
        RMSprop(learning_rate=0.001),
        Adam(learning_rate=0.001),
    ]

    base_model = MLP(seed=42)
    initial_params = base_model.copy_params()
    results = {}
    models = {}

    for opt in optimizers:
        print(f"\n{'='*50}\n{opt.get_name()}\n{'='*50}")
        model = MLP(seed=42)
        model.set_params([p.copy() for p in initial_params])
        opt.reset()
        # Train using VALIDATION set for monitoring
        results[opt.get_name()] = train(model, opt, X_train, y_train,
                                         X_val, y_val, epochs=20)
        models[opt.get_name()] = model

    # Final test evaluation (done only once!)
    print(f"\n{'='*50}\nFINAL TEST EVALUATION\n{'='*50}")
    for name, model in models.items():
        test_acc = compute_accuracy(model, X_test, y_test)
        results[name]['final_test_acc'] = test_acc
        print(f"{name.split('(')[0]:<20} Test: {test_acc*100:.2f}%")

    # Plot results (using validation accuracy, not test)
    fig, axes = plt.subplots(1, 2, figsize=(12, 4))
    colors = ['#1f77b4', '#ff7f0e', '#2ca02c', '#d62728']

    for (name, hist), color in zip(results.items(), colors):
        axes[0].plot(hist['train_loss'], label=name.split('(')[0], color=color)
        axes[1].plot([a*100 for a in hist['val_acc']], label=name.split('(')[0],
                    color=color)

    axes[0].set_xlabel('Epoch'); axes[0].set_ylabel('Loss')
    axes[0].set_title('Training Loss'); axes[0].legend(); axes[0].set_yscale('log')
    axes[1].set_xlabel('Epoch'); axes[1].set_ylabel('Accuracy (%)')
    axes[1].set_title('Validation Accuracy'); axes[1].legend()
    axes[1].axhline(y=95, color='gray', linestyle='--', alpha=0.5)

    plt.tight_layout()
    plt.savefig('optimizer_showdown.png', dpi=150)
    plt.show()

    return results


if __name__ == "__main__":
    results = run_showdown()
\end{pythoncode}

%------------------------------------------------------------------------------
\section{Summary}
\label{sec:capstone-summary}

In this capstone project, you:

\begin{itemize}
    \item \textbf{Implemented four optimizers from scratch}: SGD, Momentum, RMSprop, and Adam. You now understand what happens inside these algorithms, not just how to call them.

    \item \textbf{Built a complete training pipeline}: Forward pass, backward pass, mini-batch training, and metric tracking. These components appear in every deep learning project.

    \item \textbf{Used proper evaluation methodology}: Train/validation/test splits ensure unbiased comparisons. The validation set guides optimizer selection during training; the test set provides a single, final estimate of generalization performance.

    \item \textbf{Conducted systematic experiments}: Fair comparison requires controlling for random initialization, data ordering, and hyperparameters.

    \item \textbf{Analyzed results quantitatively}: Loss curves, accuracy trajectories, and learning rate sensitivity reveal optimizer characteristics that matter in practice.

    \item \textbf{Developed practical intuition}: Adam is the safe default; SGD+Momentum may squeeze out extra accuracy with tuning; RMSprop excels for recurrent networks.
\end{itemize}

\begin{keyconcept}
\textbf{What You Can Now Do:}
\begin{itemize}
    \item Choose the right optimizer for a new problem
    \item Debug training issues by understanding optimizer internals
    \item Implement custom optimizers when needed
    \item Run fair comparisons for research or hyperparameter tuning
\end{itemize}
\end{keyconcept}

The optimizers you've implemented here form the foundation of modern deep learning training. Variants of these algorithms---particularly Adam and AdamW---are used in many state-of-the-art systems. The math scales from MNIST to trillion-parameter models; while production systems often include additional engineering optimizations, the core principles remain the same.

\end{document}
